
\(x\in P_\bullet\), on ait \(x^2=0\)), alors il en est de même pour \(\bigoplus_{i\ge0}\Tor_i^A(B,C)\).

Cette dernière remarque s'applique au cas où \(B\) est le quotient de \(A\) par un idéal régulier, car on peut alors prendre pour \(P_\bullet\) le complexe de Koszul associé à un homomorphisme régulier définissant l'idéal de \(B\).

\subsection*{2.3.}
La structure de \(B\otimes_A C\)-algèbre définie en 2.1 sur \(\bigoplus_{i\ge0}\Tor_i^A(B,C)\) est fonctorielle en \(B\) et \(C\). On peut alors globaliser les constructions précédentes: soient \(X\) un schéma, \(X'\), \(X''\) deux \(X\)-schémas; alors on peut localement sur \(X'\times_X X''\) définir les morphismes (2.1.1), ils se recollent par fonctorialité, et définissent sur \(\bigoplus_{i\ge0}\underline{\Tor}_i^{\underline{O}_X}(\underline{O}_{X'},\underline{O}_{X''})\) une structure de \(\underline{O}_{X'\times_X X''}\)-Algèbre graduée associative.

\textbf{Lemme 2.4.} \uline{Soient} \(X',X''\) \uline{deux} \(X\)-\uline{schémas}, \uline{et} \(\underline{M}\) \uline{un} \(\underline{O}_{X'}\)-\uline{Module plat}. \uline{Alors pour tout} \(i\ge0\), \uline{on a un isomorphisme canonique} :
\begin{equation}\notag
\underline{\Tor}_i^{\underline{O}_X}(\underline{M},\underline{O}_{X''})\simeq
\underline{M}\otimes_{\underline{O}_{X'}}\underline{\Tor}_i^{\underline{O}_X}(\underline{O}_{X'},\underline{O}_{X''}).
\end{equation}

Il suffit de définir un isomorphisme fonctoriel dans le cas affine; on pourra alors par recollement le définir globalement dans le cas général.

Soient alors \(X=\Spec(A)\), \(X'=\Spec(A')\), \(X''=\Spec(A'')\); \(\underline{M}\) est défini par un \(A'\)-module plat \(M\). Soit \(P_\bullet\) une résolution projective de \(A''\) sur \(A\); les \(\Tor_i^A(A',A'')\) sont les groupes d'homologie de \(A'\otimes_A P_\bullet\). Comme \(M\) est plat sur \(A'\), on a un isomorphisme canonique : \(M\otimes_{A'}\Tor_i^A(A',A'')\xrightarrow{\sim}H_i\bigl(M\otimes_{A'}(A'\otimes_A P_\bullet)\bigr)\). Par ailleurs, l'isomorphisme canonique : \(M\otimes_{A'}(A'\otimes_A P_\bullet)\xrightarrow{\sim}M\otimes_A P_\bullet\) donne un isomorphisme canonique : \(H_i\bigl(M\otimes_{A'}(A'\otimes_A P_\bullet)\bigr)\xrightarrow{\sim}\Tor_i^A(M,A'')\), d'où le résultat.

\textbf{Proposition 2.5.} \uline{Soient} \(X\) \uline{un schéma}, \(Y\) \uline{et} \(Z\) \uline{deux sous-schémas fermés de} \(X\), \(T=Y\times_X Z\). \uline{On suppose que pour tout} \(x\in T\), \uline{il existe un voisinage ouvert}
